% Fragment partage : inclus par compte-business.tex (dossier autonome)
% et par dossier-conception-alanya.tex (dossier client assemble).
% Ne pas y mettre de preambule ni de \part : c'est l'appelant qui les pose.

\chapter{Ce qu'un commerçant peut présenter}

\section{La fiche}

Un compte business dispose d'une fiche étendue, absente d'un compte personnel~:

\begin{center}
\begin{tabularx}{\linewidth}{@{}p{4.4cm}X@{}}
\toprule
\thd{Identité} & Raison sociale, description, logo, image de couverture. \\
\thd{Catégorie} & Une catégorie principale, choisie dans une liste courte. \\
\thd{Mots-clés} & Saisis librement par le commerçant, pour être trouvé. \\
\thd{Lieu} & Adresse, position sur la carte, distance affichée au visiteur. \\
\thd{Horaires} & Jours et créneaux d'ouverture, exceptions. \\
\thd{Liens} & Site web, réseaux sociaux, dans l'ordre choisi. \\
\thd{Médias} & Images et vidéos de présentation. \\
\thd{Moyens de paiement} & Ce que l'établissement accepte sur place. \\
\thd{Catalogue} & Produits ou services, avec prix et disponibilité. \\
\bottomrule
\end{tabularx}
\end{center}

Trois de ces rubriques méritent d'être détaillées, parce qu'elles ne sont pas de simples
champs de formulaire.

\section{Le lieu et la géolocalisation}

C'est la rubrique dont dépend tout l'annuaire, et la seule qui parte réellement de zéro ---
Alanya ne stocke aujourd'hui aucune donnée géographique.

\begin{center}
\begin{tabularx}{\linewidth}{@{}p{4.4cm}X@{}}
\toprule
\thd{Adresse écrite} & Ce que le visiteur lit et recopie.
\guillemotleft~Avenue Kennedy, Yaoundé~\guillemotright. Texte libre~: il n'existe pas de
référentiel d'adresses fiable sur le marché visé. \\
\thd{Position} & Un point sur la carte, posé par le commerçant lui-même. C'est
\emph{elle} qui sert à la recherche par proximité, pas l'adresse écrite. \\
\thd{Distance affichée} & Calculée à la volée entre le visiteur et le commerce.
\guillemotleft~à 1,2~km de vous~\guillemotright. \\
\bottomrule
\end{tabularx}
\end{center}

\begin{note}[Comment le commerçant pose son point]
Deux gestes seulement~: \guillemotleft~utiliser ma position actuelle~\guillemotright{}
quand il est sur place, ou déplacer un repère sur une carte. Pas de géocodage d'adresse ---
il donnerait des résultats médiocres et masquerait les erreurs derrière une fausse
précision. Le commerçant sait où est sa boutique~; on lui demande de le montrer.
\end{note}

\begin{constat}[Aucune dépendance à ajouter]
\code{geolocator} et \code{flutter\_map} sont \textbf{déjà} dans \code{pubspec.yaml} --- ils
servent au partage de position dans les conversations
(\code{lib/screens/chats/location\_picker\_screen.dart}). Le sélecteur de position du
commerçant peut réutiliser cet écran.
\end{constat}

\section{Médias de présentation}

Images \emph{et} vidéos. C'est ce qui fait la différence entre une fiche remplie et une
fiche convaincante --- et c'est aussi la rubrique la plus coûteuse à exploiter.

\begin{center}
\begin{tabularx}{\linewidth}{@{}p{4.4cm}X@{}}
\toprule
\thd{Couverture} & Une image en tête de fiche, format paysage. \\
\thd{Galerie} & Photos de l'établissement, des produits, de l'équipe. Ordre choisi par le
commerçant. \\
\thd{Vidéos} & Présentation de l'activité. Avec vignette, durée, et lecture au clic. \\
\thd{Logo} & Réutilise l'avatar du compte --- pas de champ supplémentaire. \\
\bottomrule
\end{tabularx}
\end{center}

\begin{alerte}[La vidéo est un chantier, pas un champ]
Une vidéo téléversée depuis un téléphone pèse couramment plusieurs dizaines de mégaoctets.
La servir telle quelle à chaque ouverture de fiche est intenable. Il faut une vignette, une
limite de durée, et un transcodage --- lequel suppose la file de jobs conçue à l'étape~3.
Voir le chapitre~\ref{ch:medias}.
\end{alerte}

\section{Liens et présence en ligne}

Site web, Facebook, Instagram, WhatsApp, TikTok\ldots{} Le commerçant en ajoute autant qu'il
veut et choisit leur ordre d'affichage.

\begin{note}[Pas une colonne par réseau]
Le point que le document d'amorce soulignait à juste titre. Une table
\code{business\_link (type, url, ordre)} permet d'ajouter le réseau à la mode de l'an
prochain sans migration. Une colonne \code{facebook\_url} ne le permettrait pas.
\end{note}

\section{Les quatre écrans}

\begin{center}
  \imageOuCadre{maquettes/comptes-business-ecrans.png}{\linewidth}%
    {Fiche business, catalogue, éditeur d'horaires, annuaire\\
     \texttt{maquettes/comptes-business-ecrans.png}}
\end{center}

\begin{note}[Maquette interactive]
\href{https://claude.ai/code/artifact/37b8e6c7-1235-42df-aabc-d66b8f26ff67}%
{\texttt{claude.ai/code/artifact/37b8e6c7}}\\
Source versionnée~: \code{docs/architecture/maquettes/comptes-business.html}
\end{note}

\section{Le badge fait son travail}

Dans l'annuaire, trois commerces se distinguent d'un coup d'{\oe}il~: deux portent la
pastille pleine \badgePastillePanier[1] --- vérifiés par Alanya --- et un porte le panier nu
\badgePanier[1], c'est-à-dire un commerce qui s'est déclaré tel sans qu'on l'ait contrôlé.

C'est exactement la distinction que le socle devait rendre lisible à 13~pixels, et c'est
dans l'annuaire qu'elle compte le plus~: un visiteur qui cherche une pharmacie de garde
prend une décision sur cette base.

\chapter{La catégorie et les mots-clés}

\section{Une liste courte, pas une arborescence}

Le choix retenu~: \textbf{une catégorie principale} dans une liste courte, plus des
\textbf{mots-clés libres} saisis par le commerçant.

\begin{center}
\begin{tabularx}{\linewidth}{@{}p{4.4cm}X@{}}
\toprule
\thd{Pourquoi pas deux niveaux} & Un sélecteur à deux étages est pénible sur un téléphone,
et il faudrait arrêter dès aujourd'hui une liste longue qu'on ne pourra plus changer. \\
\thd{Pourquoi des mots-clés} & Ils absorbent tout ce que la liste ne prévoit pas ---
\guillemotleft~garde de nuit~\guillemotright, \guillemotleft~livraison~\guillemotright,
\guillemotleft~sur mesure~\guillemotright{} --- sans figer quoi que ce soit. \\
\bottomrule
\end{tabularx}
\end{center}

\begin{alerte}[Le revers, à assumer]
La qualité de la recherche dépendra de ce que les commerçants écrivent. Deux garde-fous
tempèrent~: les mots-clés sont \emph{normalisés} (minuscules, accents repliés) et le champ
propose les mots déjà utilisés par d'autres commerces de la même catégorie. Un mot-clé
partagé vaut mieux que dix orthographes du même mot.
\end{alerte}

\section{Une taxonomie est difficile à changer}

C'est le point que le document d'amorce soulignait à juste titre~: une fois des milliers de
fiches rangées, renommer ou fusionner une catégorie devient une opération de reprise de
données. D'où la liste courte --- vingt à trente entrées --- et la table de référence
éditable plutôt que des valeurs figées dans le code.

\chapter{Les horaires, l'écran le plus difficile}

Ce n'est pas un champ de formulaire. Six cas doivent tenir sur un écran de téléphone, et
cinq d'entre eux restent invisibles tant qu'on n'a pas essayé.

\begin{center}
\begin{tabularx}{\linewidth}{@{}p{4.1cm}X@{}}
\toprule
\thd{Multi-créneaux} & Ouvert le matin, fermé à midi, rouvert l'après-midi. Deux plages pour
un seul jour --- le cas le plus courant, et celui qu'un champ \guillemotleft~de\ldots{}
à\ldots{}~\guillemotright{} ne sait pas exprimer. \\
\thd{24\,h/24} & N'est pas \guillemotleft~00:00 $\rightarrow$ 00:00~\guillemotright. C'est un
état distinct, sans quoi le calcul \guillemotleft~ouvert maintenant~\guillemotright{} se
trompe. \\
\thd{Passage de minuit} & Un bar ouvert de 20\,h à 2\,h~: l'heure de fin est
\emph{avant} l'heure de début. \\
\thd{Copie d'un jour} & Saisir sept fois les mêmes horaires fait abandonner.
\guillemotleft~Copier vers~\guillemotright{} n'est pas un confort, c'est ce qui décide si la
fiche sera remplie. \\
\thd{Exceptions} & Jours fériés, congés, fermeture ponctuelle. Une liste à part, pas un
champ de plus. \\
\thd{Fuseau horaire} & \guillemotleft~Ouvert maintenant~\guillemotright{} se calcule dans le
fuseau du commerce, pas dans celui du visiteur. \\
\bottomrule
\end{tabularx}
\end{center}

\begin{constat}[Le fuseau existe déjà]
La table \code{pays} porte \code{timeZone} et \code{decalageHoraire}
(\code{migrations/001\_initial\_schema.sql}), et \code{users.idPays} les relie. Rien à
créer --- mais tout à ne pas oublier.
\end{constat}

\begin{note}[À budgéter comme un écran]
Cet éditeur mérite sa propre estimation, au même titre qu'un écran de conversation. Le
traiter comme un champ du formulaire de profil garantit qu'il sera bâclé, puis refait.
\end{note}

\chapter{Le tableau de bord du commerçant}

\section{Ce qu'il montre}

Puisque les transactions se font hors application, le tableau de bord ne peut pas parler
d'argent. Il parle de ce qu'Alanya \textbf{observe réellement}~: l'audience de la fiche.

\begin{itemize}[leftmargin=1.4em]
  \item nombre de vues du profil, et son évolution~;
  \item provenance des visites --- annuaire, recherche, conversation, lien partagé~;
  \item répartition géographique des visiteurs~;
  \item conversations entamées depuis la fiche.
\end{itemize}

\begin{note}[Pourquoi pas de commandes déclarées]
On aurait pu laisser le commerçant saisir ses ventes pour obtenir un
\guillemotleft~suivi des transactions~\guillemotright. L'expérience de ce genre de saisie
manuelle est constante~: elle est abandonnée après quelques semaines, et le tableau de bord
devient faux --- pire que vide. Écarté volontairement.
\end{note}

\section{Ce que cela implique}

Compter des vues suppose de les enregistrer. C'est un journal d'événements, avec le volume
que cela suppose --- voir le chapitre technique~8.
